
\documentclass{article}
\usepackage{colortbl}
\usepackage{makecell}
\usepackage{multirow}
\usepackage{supertabular}

\begin{document}

\newcounter{utterance}

\centering \large Interaction Transcript for game `cladder', experiment `full\_v1.5\_default', episode 5168 with refineonly{-}llama{-}3.1{-}8b{-}instruct.
\vspace{24pt}

{ \footnotesize  \setcounter{utterance}{1}
\setlength{\tabcolsep}{0pt}
\begin{supertabular}{c@{$\;$}|p{.15\linewidth}@{}p{.15\linewidth}p{.15\linewidth}p{.15\linewidth}p{.15\linewidth}p{.15\linewidth}}
    \# & \multicolumn{2}{c}{Player} && \multicolumn{2}{c}{Game Master} \\
    \hline

    \theutterance \stepcounter{utterance}  
    & & & \multicolumn{4}{p{0.6\linewidth}}{
        \cellcolor[rgb]{0.9,0.9,0.9}{
            \makecell[{{p{\linewidth}}}]{
                \texttt{\tiny{[P1$\langle$GM]}}
                \texttt{You are an expert in causal inference. The following question is not a typical commonsense query, but rather a meticulously designed question created by a professor specializing in causal inference, intended to assess the students' mastery of the course content.} \\
\texttt{Imagine a self{-}contained, hypothetical world with only the following conditions, and without any unmentioned factors or causal relationships: Talent has a direct effect on elite institution admission status. Brown eyes has a direct effect on elite institution admission status. The overall probability of talent is 97\%. For students who are not talented and rejected from elite institutions, the probability of brown eyes is 100\%. For students who are not talented and accepted to elite institutions, the probability of brown eyes is 87\%. For students who are talented and rejected from elite institutions, the probability of brown eyes is 97\%. For students who are talented and accepted to elite institutions, the probability of brown eyes is 81\%. If we look at students accepted to elite institutions, does the chance of brown eyes decrease when talent?} \\
\texttt{Start your answer with 'Yes' or 'No', followed by additional reasoning or evidence to support your explanation.} \\
            }
        }
    }
     \\ \\

    \theutterance \stepcounter{utterance}  
    & & \multicolumn{4}{p{0.6\linewidth}}{
        \cellcolor[rgb]{0.95,0.95,0.95}{
            \makecell[{{p{\linewidth}}}]{
                \texttt{\tiny{[P1$\leftrightarrow$P]}}
                \texttt{{[}DRAFT{]}} \\
\texttt{No.} \\
\\ 
\texttt{To answer this question, we need to apply the concept of conditional probability and the causal relationships given in the problem. We are interested in the probability of brown eyes among students who are accepted to elite institutions, given that they are talented.} \\
\\ 
\texttt{Let's denote the following events:} \\
\texttt{{-} T: Talent} \\
\texttt{{-} B: Brown eyes} \\
\texttt{{-} A: Accepted to elite institutions} \\
\texttt{{-} R: Rejected from elite institutions} \\
\\ 
\texttt{We are given the following conditional probabilities:} \\
\texttt{{-} P(T) = 0.97 (probability of talent)} \\
\texttt{{-} P(B|R, ¬T) = 1 (probability of brown eyes among non{-}talented students who are rejected)} \\
\texttt{{-} P(B|R, T) = 0.97 (probability of brown eyes among talented students who are rejected)} \\
\texttt{{-} P(B|A, ¬T) = 0.87 (probability of brown eyes among non{-}talented students who are accepted)} \\
\texttt{{-} P(B|A, T) = 0.81 (probability of brown eyes among talented students who are accepted)} \\
\\ 
\texttt{We want to find P(B|A, T), the probability of brown eyes among students who are accepted to elite institutions, given that they are talented. To do this, we can use the law of total probability:} \\
\\ 
\texttt{P(B|A, T) = P(B|A, T, T) * P(T|A, T) + P(B|A, T, ¬T) * P(¬T|A, T)} \\
\\ 
\texttt{However, we don't have direct information about P(T|A, T) and P(¬T|A, T). We can use the fact that P(T) = 0.97 and the definition of conditional probability to rewrite these expressions:} \\
\\ 
\texttt{P(T|A, T) = 1 (since T is given)} \\
\texttt{P(¬T|A, T) = 0 (since T is given)} \\
\\ 
\texttt{Now, we need to find P(B|A, T, T) and P(B|A, T, ¬T). We are given P(B|A, T) = 0.81 and P(B|A, ¬T) = 0.87. We can use the law of total probability again to express P(B|A, T) in terms of P(B|A, T, T) and P(B|A, T, ¬T):} \\
\\ 
\texttt{P(B|A, T) = P(B|A, T, T) * P(T|A, T) + P(B|A, T, ¬T) * P(¬T|A, T)} \\
\texttt{= P(B|A, T, T) * 1 + P(B|A, T, ¬T) * 0} \\
\texttt{= P(B|A, T, T)} \\
\\ 
\texttt{Similarly, we can express P(B|A, ¬T) in terms of P(B|A, ¬T, T) and P(B|A, ¬T} \\
            }
        }
    }
    & \\ \\

    \theutterance \stepcounter{utterance}  
    & \multicolumn{4}{p{0.6\linewidth}}{
        \cellcolor[rgb]{0.561,0.737,0.561}{
            \makecell[{{p{\linewidth}}}]{
                \texttt{\tiny{[P1$\rangle$GM]}}
                \texttt{No.} \\
\\ 
\texttt{To answer this question, we need to apply the concept of conditional probability and the causal relationships given in the problem. We are interested in the probability of brown eyes among students who are accepted to elite institutions, given that they are talented.} \\
\\ 
\texttt{Let's denote the following events:} \\
\texttt{{-} T: Talent} \\
\texttt{{-} B: Brown eyes} \\
\texttt{{-} A: Accepted to elite institutions} \\
\texttt{{-} R: Rejected from elite institutions} \\
\\ 
\texttt{We are given the following conditional probabilities:} \\
\texttt{{-} P(T) = 0.97 (probability of talent)} \\
\texttt{{-} P(B|R, ¬T) = 1 (probability of brown eyes among non{-}talented students who are rejected)} \\
\texttt{{-} P(B|R, T) = 0.97 (probability of brown eyes among talented students who are rejected)} \\
\texttt{{-} P(B|A, ¬T) = 0.87 (probability of brown eyes among non{-}talented students who are accepted)} \\
\texttt{{-} P(B|A, T) = 0.81 (probability of brown eyes among talented students who are accepted)} \\
\\ 
\texttt{We want to find P(B|A, T), the probability of brown eyes among students who are accepted to elite institutions, given that they are talented. To do this, we can use the law of total probability:} \\
\\ 
\texttt{P(B|A, T) = P(B|A, T, T) * P(T|A, T) + P(B|A, T, ¬T) * P(¬T|A, T)} \\
\\ 
\texttt{However, we don't have direct information about P(T|A, T) and P(¬T|A, T). We can use the fact that P(T) = 0.97 and the definition of conditional probability to rewrite these expressions:} \\
\\ 
\texttt{P(T|A, T) = 1 (since T is given)} \\
\texttt{P(¬T|A, T) = 0 (since T is given)} \\
\\ 
\texttt{Now, we need to find P(B|A, T, T) and P(B|A, T, ¬T). We are given P(B|A, T) = 0.81 and P(B|A, ¬T) = 0.87. We can use the law of total probability again to express P(B|A, T) in terms of P(B|A, T, T) and P(B|A, T, ¬T):} \\
\\ 
\texttt{P(B|A, T) = P(B|A, T, T) * P(T|A, T) + P(B|A, T, ¬T) * P(¬T|A, T)} \\
\texttt{= P(B|A, T, T) * 1 + P(B|A, T, ¬T) * 0} \\
\texttt{= P(B|A, T, T)} \\
\\ 
\texttt{Similarly, we can express P(B|A, ¬T) in terms of P(B|A, ¬T, T) and P(B|A, ¬T} \\
            }
        }
    }
    & & \\ \\

    \theutterance \stepcounter{utterance}  
    & & & \multicolumn{2}{p{0.3\linewidth}}{
        \cellcolor[rgb]{0.9,0.9,0.9}{
            \makecell[{{p{\linewidth}}}]{
                \texttt{\tiny{[GM$|$GM]}}
                \texttt{No} \\
            }
        }
    }
    & & \\ \\

    \theutterance \stepcounter{utterance}  
    & & & \multicolumn{2}{p{0.3\linewidth}}{
        \cellcolor[rgb]{0.9,0.9,0.9}{
            \makecell[{{p{\linewidth}}}]{
                \texttt{\tiny{[GM$|$GM]}}
                \texttt{yes} \\
            }
        }
    }
    & & \\ \\

    \theutterance \stepcounter{utterance}  
    & & & \multicolumn{2}{p{0.3\linewidth}}{
        \cellcolor[rgb]{0.9,0.9,0.9}{
            \makecell[{{p{\linewidth}}}]{
                \texttt{\tiny{[GM$|$GM]}}
                \texttt{game\_result = LOSE} \\
            }
        }
    }
    & & \\ \\

\end{supertabular}
}

\end{document}
